\centerline{\bf \Huge Прогрессии}

\begin{flushright}
        \scriptsize{}
\end{flushright}

\problem{1}
Дана арифметическая прогрессия. Сумма первых её 10 членов равна 60 , а сумма первых 20 её членов равна 320. Чему может быть равен 15-й член этой прогрессии?

\problem{2}
В арифметической прогрессии сумма ее $m$ первых членов равна сумме $n$ первых членов ($m$ не равно $n$). Докажите, что в этом случае сумма ее первых $m+n$ членов равна нулю.

\problem{3}
Две арифметические прогрессии содержат по 2015 членов каждая. Отношение последнего члена первой прогрессии к первому члену второй равно отношению последнего члена второй прогрессии к первому члену первой и равно 4 . Отношение суммы всех членов первой прогрессии к сумме всех членов второй равно 2. Найдите отношение разностей этих прогрессий и приведите пример таких прогрессий.

\problem{4}
Найти восемнадцатый член арифметической прогрессии, если первый и одиннадцатый её члены - натуральные числа, а сумма первых четырнадцати членов равна 77.

\problem{5}
Найти первый член целочисленной арифметической прогрессии, у которой сумма первых шести членов отличается от суммы следующих шести членов менее чем на 450, а сумма первых пяти членов превышает более чем на 5 сумму любого другого набора различных членов этой прогрессии.

\problem{6}
В бесконечной числовой последовательности $x_1, x_2, \ldots, x_n, \ldots$ не все члены равны между собой. Для всех $n \geqslant 2$ выполняется равенство

$$
x_n=\frac{x_{n-1}+x_n+x_{n+1}}{3} .
$$


Найдите отношение $\frac{x_{2012}-x_{1006}}{x_{1006}-x_{503}}$.